% ============================================================================
% Radiology: Artificial Intelligence: Original Research
% Written by scripts/configure.py from journals/requirements.json.
%
% A community template: not made or endorsed by the RSNA. The lines below are what was read on
% 2026-11-21, in plain words. They are not the journal's own words, and nobody has compared them
% with the live page. Check the journal's instructions before you submit:
%   https://pubs.rsna.org/page/ai/author-instructions
%
% How sure each line is, in [brackets]:
%   A: read on the official site itself on the date given
%   B: read in a full copy of the official page (the browser print of 2025-09-18), or the same value
%      in two or more separate answers of the search tool about the official page, or one such
%      answer in the page's own words that the journal's printed articles agree with. Agreement
%      between answers shows that the copy was read steadily, not that the copy is current
%   C: one answer of the search tool only, or an answer that was ambiguous; confirm on the live page
%      first
%   P: measured from the journal's printed articles; shows what is printed after editing, not what
%      the instructions ask
%
% A number is an upper limit unless its line says otherwise.
%
% This article type:
%   By invitation: no [B]
%   Body words: 3000 [B]
%   What the body count covers: Introduction to Discussion [B]
%   Words by section: Introduction: 400; Materials and Methods: 800; Results: 1000; Discussion: 800
%       [C]
%   Abstract: structured [B]
%   Abstract words: 250 [B]
%   Abstract headings: Purpose, Materials and Methods, Results, Conclusion [P]
%   Summary statement: required: yes; characters at most: 255; sentences: 1 (the manuscript-type
%       entries say 1-2 sentences; the title-page entry and the checklist say a single sentence of
%       no more than 255 characters, in boldface) [B]
%   Key points box: Key Points [B]
%   Key points, how many: 3 [B]
%   Key points, words: not read; check the journal's page
%   Body headings: Introduction, Materials and Methods, Results, Discussion [P]
%   References: 35 [B]
%   Figures: 6 [B]
%   Tables: 4 [B]
%   Reporting checklist: when: at first submission of Original Research and Technical Development
%       papers that involve human subjects; not for animal studies; mark items that do not apply
%       N/A; CLAIM: for AI manuscripts in medical imaging; STARD: diagnostic accuracy; CONSORT:
%       randomized controlled trials; PRISMA: meta-analyses; STROBE: observational studies (cohort,
%       case-control, cross-sectional) [B]
%   Also: systematic reviews and meta-analyses are submitted as Original Research [B]
%
% This journal, every article type:
%   Peer review: double-anonymized [C]
%   Title words: not read; check the journal's page
%   Title pages: full title page: a separate file: all authors, institutions and academic degrees;
%       affiliations by superscript numbers in the order of authorship; funding and the data sharing
%       statement belong here (RSNA policies page); abbreviated title page: anonymized, the first
%       page of the manuscript file: full title, article type, five keywords, abbreviations (no more
%       than 10), the summary statement in boldface, the Key Points [B]
%   Anonymisation: the abbreviated title page and the acknowledgments are anonymized; no author
%       institution and no funding information in the manuscript file; authors' initials are
%       anonymized where the authors are the readers [B]
%   Abbreviations: at most: 10; listed on: abbreviated title page; not in: the Key Points (obvious
%       ones like CT and MRI are fine) [B]
%   Order of the manuscript file: abbreviated title page (anonymized), abstract, introduction,
%       materials and methods, results, discussion, acknowledgments (anonymized), references, figure
%       legends, tables (one per page) [C]
%   Separate files: full title page; cover letter (sole submission to this journal; conflicts of
%       interest; overlap with earlier publications; an explanation where there are two first
%       authors); reporting checklist; supplemental material in one Word file; figures [C]
%   File type: Microsoft Word (.docx) or PDF, per the checklist page; the What's New entry of
%       January 2026 says the text is uploaded as .docx and .doc is no longer accepted. Whether a
%       PDF made from LaTeX is still taken for the first submission was not confirmed. [C]
%   Font: a basic font such as Arial, Calibri or Times New Roman, 11 points [C]
%   Spacing: double-spaced, adequate margins, not right-justified [C]
%   Page numbers: none: the system numbers the pages on conversion to PDF [C]
%   Line numbers: not read; check the journal's page
%   Authors: one corresponding or senior author; at most two first authors; more than 40 authors
%       needs the Editor's consent; CRediT taxonomy for certain manuscript types (since October
%       2025) [C]
%
% Statements, for every article type of this journal:
%   Funding: on the full title page only; never in the anonymized manuscript file [C]
%   Disclosures: conflicts of interest are stated in the cover letter [C]
%   Data sharing: a data sharing statement on the full title page (RSNA policy for all its journals)
%       [C]
%   Code availability: not read; check the journal's page
%   Use of AI tools in writing: disclosed at submission; where used, described in both the cover
%       letter and the manuscript; an AI tool cannot be an author [B]
%   Reporting checklists: when: at first submission of Original Research and Technical Development
%       papers that involve human subjects; not for animal studies; mark items that do not apply
%       N/A; CLAIM: for AI manuscripts in medical imaging; STARD: diagnostic accuracy; CONSORT:
%       randomized controlled trials; PRISMA: meta-analyses; STROBE: observational studies (cohort,
%       case-control, cross-sectional) [B]
%   Trial registration: not read; check the journal's page
%   Ethics approval: not read; check the journal's page
%
% The writing notes under the four headings of a research article are kept from the template this
% one descends from. They were written for Radiology; the limits in them were replaced by the lines
% read for this journal.
% ============================================================================

% Use basic fonts such as Arial, Calibri, or Times New Roman at 11 points.
\documentclass[11pt]{article}
% Use basic fonts such as Arial, Calibri, or Times New Roman at 11 points.
\usepackage{fontspec}
\setmainfont{Times New Roman}

% Enable html conversion through lwarp
\usepackage[mathjax]{lwarp}

\makeatletter
\@ifpackageloaded{lwarp}{%
  % lwarp is loaded
  % Instruct lwarp to help generate html output for MS Word reading
  \booltrue{FormatWP}
  
  % Specify font and other styles for html output
  \CSSFilename{rsnastyle.css}
}{%
  % lwarp is NOT loaded
}
\makeatother

% Hack to make the bibliography into a numbered list in MS Word
% See https://github.com/bdtc/lwarp/discussions/38
\makeatletter

\newcommand*{\LWR@enumerateitem}[1][]{%
    \LWR@stoppars%
    \LWR@startnewdepth{listitem}%
    \LWR@htmltag{li}%
    \LWR@orignewline%
    \LWR@startpars%
    \LWR@ensuredoingapar%
    \LWR@origitem[#1]%
}

\def\LWR@enumeratestart{%
    \renewcommand*{\LWR@printcloselist}{\LWR@printcloseenumerate}%
    \renewcommand*{\LWR@printopenlist}{%
        ol % space
        class=\textquotedbl{}enumerate\textquotedbl{} % space
        style=\textquotedbl\LWR@print@mbox{list-style-type:decimal}\textquotedbl{}%
    }%
    \LetLtxMacro\item\LWR@enumerateitem%
    \LWR@nulllistfills%
}

\makeatother

\usepackage{graphicx} % Required for inserting images

% Page size
\usepackage[paper=a4paper,top=2.5cm,bottom=2.5cm,left=3cm,right=3cm]{geometry}

% Remove page numbers
\pagestyle{empty}

% Double line spacing
\usepackage{setspace}
\doublespacing

% Author info
\usepackage{authblk}
\date{}

% Anonymous title style
\makeatletter
\renewcommand{\maketitle}{
\begin{center}
{\LARGE\textbf{\@title}}
\end{center}
}
\makeatother


% Push the figures to the end
\usepackage[nolists,tablesfirst,nomarkers]{endfloat}

% Customize List of Figures and Tables
\usepackage{tocloft}
\renewcommand{\cftfigpresnum}{Fig.~}   % Figure number prefix
%\renewcommand{\cftfigaftersnum}{} % Figure number suffix
\setlength{\cftfigindent}{0pt}    % left align everything
\setlength{\cftfignumwidth}{1.2cm}  % space reserved for numbers
\cftpagenumbersoff{figure}

\renewcommand{\cfttabpresnum}{Table~}
\setlength{\cfttabindent}{0pt}    % left align everything
\setlength{\cfttabnumwidth}{1.5cm}  % space reserved for numbers
\cftpagenumbersoff{table}

% Remove indentation of paragraphs
\usepackage{parskip}
\setlength{\parindent}{0pt} % No indentation
%\setlength{\parskip}{1em}   % Space between paragraphs

% Number references consecutively in the order in which they are first mentioned in the manuscript.
% Verify reference order at the time of revision.
% The abbreviations used for periodicals cited in the references should follow the style of the National Library of Medicine.
% It is the responsibility of the author(s) to ensure the accuracy of all references.
% Citation format:
% 1. Tyler LN, Harville TO, Blackall DP, Callegari PE, Keenan GF, Elliott M. Multiple alloantibodies after transfusion in an infant treated with infliximab. N Engl J Med 2007;357:2092-2093. doi: 10.1056/NEJMc070741
\usepackage[style=nejm,natbib=true]{biblatex}

% Use a basic enumerate for the bibliography for compatibility with MS Word
\defbibenvironment{simplebib}
  {\enumerate}
  {\endenumerate}
  {\item}  

% Handy commands
\newcommand{\figref}[1]{\figurename~\ref{#1}}
\newcommand{\tabref}[1]{\tablename~\ref{#1}}

% use normal font for URLs
\urlstyle{same}

% Don't use hyperref for simplicity of MS Word conversion
%\usepackage[hidelinks]{hyperref}

% If you use hyperref and glossaries, you must load hyperref first (although, in general, hyperref should be loaded after other packages).
\usepackage[acronym]{glossaries} % for abbreviations
\renewcommand*{\glossarysection}[2][]{} % removes heading completely
\makeglossaries

\addbibresource{sample_bibliography.bib}
\addbibresource{xampl.bib}

% For lorem ipsum only, can be removed
\usepackage{mwe}
\usepackage{blindtext}

% Title words: not read; check the journal's page
% Title of the manuscript. Include modality and disease. Be concise. If space allows, indicate study type (randomized clinical trial) and registry or clinical trial name (eg, MESA, DETERMINE)
\title{Manuscript title}

% Authors: none in this file. It is the anonymized manuscript and prints the title alone.
% The authors, their institutions, the funding and the data sharing statement go on the
% full title page, a separate file: title_page.tex.

\begin{document}

\maketitle

% ----------------------------------------------------------------------------
% The first page of the manuscript file: the abbreviated title page.
% What goes on it: anonymized, the first page of the manuscript file: full title, article type, five
% keywords, abbreviations (no more than 10), the summary statement in boldface, the Key Points [B]
% ----------------------------------------------------------------------------

\textbf{Article type}: Original Research

% Keywords: named among the items of this page (see the line above).
\textbf{Keywords}: TODO

% Summary statement: required: yes; characters at most: 255; sentences: 1 (the manuscript-type
% entries say 1-2 sentences; the title-page entry and the checklist say a single sentence of no more
% than 255 characters, in boldface) [B]
% The summary statement is in boldface. It highlights the key finding in your study.
\textbf{Summary statement: TODO}

% Key points box: Key Points [B]
% Key points, how many: 3 [B]
% Key points, words: not read; check the journal's page
% The main results/conclusions from your study, including summary data. Do NOT repeat your summary statement. Please avoid using vague language and abbreviations here. Obvious abbreviations like CT and MRI are fine.
\textbf{Key Points}:
\begin{itemize}
    \item TODO
    \item TODO
    \item TODO
\end{itemize}

% Abbreviations: at most: 10; listed on: abbreviated title page; not in: the Key Points (obvious
% ones like CT and MRI are fine) [B]
% List of all abbreviations with spelled out words
\newacronym{MRI}{MRI}{Magnetic Resonance Imaging}
\newacronym{CT}{CT}{Computed Tomography}

\textbf{List of abbreviations}:
\printglossary[
  type=\acronymtype,
  title={},      % empty title
  toctitle={},   % avoid a ToC title
  nonumberlist,  % suppress page list
  nopostdot      % remove trailing do
  ]

\clearpage
% Abstract: structured [B]
% Abstract words: 250 [B]
% Abstract headings: Purpose, Materials and Methods, Results, Conclusion [P]

\section*{Abstract}

% Purpose: In one to two sentences, succinctly state the study purpose.
Purpose: We will talk about \acrfull{MRI} or \acrshort{MRI}. \blindtext

% Materials and Methods: (1) Identify as prospective or retrospective; (2) Date range of study; (3) If animal, study, provide number and type of animals; (4) Describe various groups, including controls; (5) Describe procedures performed; (6) Specifics of evaluation should parallel the results portion of the abstract; (7) Include one sentence on statistical analyses
Materials and Methods: \blindtext

% Results: The first sentence should indicate the number, sex, and mean age ± standard deviation or median age and interquartile range (whole numbers only) of the study participants enrolled. For example, if there are 100 patients of 60 men and 40 women, then write 100 patients (mean age, 47 years +/- 10 [standard deviation]; 60 men) were evaluated. Provide sex of the larger number of patients. Then, provide the findings in parallel with the Methods and include indicators of statistical significance, actual numbers, and percentages. Report the numerator and denominator for all percentages. All p-values need corresponding comparison values. Include confidence intervals when appropriate.
Results: \blindtext

% Conclusion: Must derive directly from the results. Do not elaborate on the significance or other implications of your study. Must address the purpose of the study.
Conclusion: \blindtext

\clearpage
% ----------------------------------------------------------------------------
% The body. Body words: 3000 [B]
% What the body count covers: Introduction to Discussion [B]
% Body headings: Introduction, Materials and Methods, Results, Discussion [P]
% References: 35 [B]
% Figures: 6 [B]
% Tables: 4 [B]
% ----------------------------------------------------------------------------

\section*{Introduction}
\label{sec:intro}
% Words in this section: 400 [C]
% Typically contains background information and references that inform the reader as to why the study was performed. No extensive literature review. Brevity and focus are important. In the final paragraph, state the hypothesis (for prospective studies only) and purpose. Your purpose should be essentially the same as in the abstract. Do not provide a hypothesis for retrospective studies, meta-analysis/systematic reviews, or technical development papers.
\blindtext

\section*{Materials and Methods}
\label{sec:methods}
% Words in this section: 800 [C]
% Include all items presented in the Results section.
%
% The Materials and Methods section should include:
% A clear statement of whether the study is retrospective or prospective. If main prospective study did not include your study purpose in original design, you may classify as a secondary analysis of a prospective trial.
% For prospective studies, enrolled patients become participants. Refer to them as such.
% Date range of the study.
% Sample selection: Describe inclusion and exclusion criteria, consecutive or random selection.
% Indicate how sample size was derived.
% Number, training, and expertise of persons executing and reading the index tests. State the number of years of experience of those who performed readings or evaluations. For resident readings, indicate “in-training.”
% Provide initials for any authors executing index tests or performing readings/evaluations. If those contributors are not authors, mention in acknowledgments but obtain written permission to do so.
% What was evaluated and how was evaluation performed? Mention anonymizing of data or images used in your analysis. Indicate if studies were read independently, if consensus agreements were used, and how disagreements were handled. Single observer studies are discouraged, as are “consensus” reads.
% Description of the pre-specified primary outcome measure(s), as well as secondary and other variables.
%
% For prospective and retrospective studies using human subjects, address in the first paragraph (1) Institutional review board (IRB) approval and patient informed consent. State if consent was obtained; (2) HIPAA compliance for U.S. studies; (3) Clinical Trials Registry Number and Registry Name. Mention whether all or part of your patient population was previously reported with citation of the appropriate reference(s). At the editor’s discretion, we may request a copy/ verification of IRB/ institutional review approval of your study prior to acceptance.
%
% For animal experiments, provide a statement of approval by the institutional animal care committee or appropriate substitute (see the Human and Animal Studies section). If the portion of the animals was previously reported, please indicated the reference.
%
% Describe clearly the number (if animals) and selection of the subjects studied (patients or experimental animals, including controls). Identify any utilized instruments or drugs (including contrast) with trade names, manufacturer’s name in parentheses. Describe procedures in sufficient detail to allow others to reproduce the study. It is essential to explain the manner in which studies were evaluated: independent readings, anonymized or unanonymized to other information, time sequencing between readings of several studies of the same subject to eliminate recall bias, random ordering of studies.
%
% Give references to established methods, including statistical methods that have been published but are not well known; describe new or substantially modified methods and give reasons for using these techniques.
%
% Last Paragraph: State the statistical methods used to analyze the data. Mention the statistical software used, along with version and manufacturer's name. State the p-value used for significance. If statistical analyst is an author, then provide initials. A statement of justification of sample size and/ or power calculation is strongly recommended.
\blindtext

\begin{figure}[tbh!]
\centering
\includegraphics[width=0.8\linewidth]{example-image.pdf}
\caption{Example Image.\label{fig:example}}
\end{figure}

\blindtext

\section*{Results}
\label{sec:results}
% Words in this section: 1000 [C]
% The first paragraph should summarize the demographics of your study sample. At a minimum, indicate number of patients/participants, mean age ± standard deviation or median age and interquartile range, and number of men vs women. For manuscripts with human study participants, include a table of demographics/characteristics. This should be Table 1 of your paper. The first paragraph of Results should summarize and cite the table 1 demographics. Also, document the number of participants excluded for each exclusion criteria described in methods. Include a flow diagram as Figure 1 and cite in the Results, showing initial number of participants and those excluded for any reason.
%
% Present the results in logical sequence in the text, along with tables and illustrations. Summarize important observations with reference to appropriate tables and figures. Give results for all items mentioned in the Materials and Methods. Use Materials and Methods subheadings in the same sequence in the Results, as appropriate.
%
% State the statistical significance of the findings. Report the results of the statistical analyses for all variables collected and analyzed, not just for those which exhibited statistical or near statistical significance. Avoid hedge terms such as “trend” or “marginal” for results that are not statistically significant.
%
% Provide numerators and denominators for all percentages, including sensitivity, specificity, accuracy, positive predictive value, and negative predictive value. Regression coefficients or mean values of compared groups must be given along with p-values.
%
% P-values should be expressed to 2 digits to the right of the decimal point (regardless of whether the p-value is significant) unless p <.01 in which case the p-value should be expressed to 3 digits to the right of the decimal point. The exception to this rule is when rounding p from 3 digits to 2 digits would result in p appearing non-significant (such as p=.046). Give exact p-value unless <.001. The smallest p-value that should be expressed is <.001. The largest p-value is >.99.
\blindtext

The problem  of XYZ has been studied before \cite{Lamport:Book:1989}.
\citet{Burkhardt:JNS:2019} wrote interesting stuff.
Some text with another citation~\cite{article-minimal}.

\begin{table}[tbh!]
\caption{Table to test captions and labels.\label{tab:example}}
\centering
\begin{tabular}{|c c c c|}
 \hline
 Col1 & Col2 & Col2 & Col3 \\ [0.5ex]
 \hline\hline
 1 & 6 & 87837 & 787 \\
 2 & 7 & 78 & 5415 \\
 3 & 545 & 778 & 7507 \\
 4 & 545 & 18744 & 7560 \\
 5 & 88 & 788 & 6344 \\ [1ex]
 \hline
\end{tabular}
\end{table}

\section*{Discussion}
\label{sec:discussion}
% Words in this section: 800 [C]
% Do not repeat in detail the data given in the Results section. Do not cite tables or figures in the discussion. These should be introduced in the Materials and Methods and Results sections.
%
% The first paragraph should briefly restate the problem that you studied and the primary results. Provide a succinct one paragraph summary of your entire study: briefly restate the background for your study and why it was done. Then state your major findings. Instead of providing general statements that something was generally better or superior, provide specific metrics (key results) and p-values that support your statements.
%
% Include implications of the findings and their limitations, in particular with reference to the use of modified methods, statistical or otherwise. Relate the observations to other relevant studies. Link the conclusions with the goals of the study, but avoid unqualified statements and conclusions not supported completely by the data. Avoid claiming priority and alluding to or giving the results of ongoing investigations that have not been completed or that are not part of the study. State new hypotheses when warranted, but clearly label them as such.
%
% The second to last paragraph of this section should address study limitations. The final paragraph should summarize your results and findings; if appropriate, offer what you think should be done in the future to advance your study.
\blindtext

% ----------------------------------------------------------------------------
% Statements. Funding, disclosures and the data sharing statement are on the full title
% page (title_page.tex), which is not anonymized. The two below belong to the manuscript:
% delete the one that does not apply.
% ----------------------------------------------------------------------------

% Use of AI tools in writing: disclosed at submission; where used, described in both the cover
% letter and the manuscript; an AI tool cannot be an author [B]
\textbf{Use of AI tools in writing}: TODO (the tool, its version and maker, and what it was used for)

% Code availability: not read; check the journal's page
\textbf{Code availability}: TODO (where the code is kept, with a link)

% Reporting checklist: when: at first submission of Original Research and Technical Development
% papers that involve human subjects; not for animal studies; mark items that do not apply N/A;
% CLAIM: for AI manuscripts in medical imaging; STARD: diagnostic accuracy; CONSORT: randomized
% controlled trials; PRISMA: meta-analyses; STROBE: observational studies (cohort, case-control,
% cross-sectional) [B]
% The completed checklist is a file of its own: it is not part of this manuscript.

% Acknowledgments: in this file, before the references, anonymized [C].
\section*{Acknowledgments}
TODO

% References
\printbibliography[env=simplebib]

\clearpage
\section*{Figure Legend(s)}
\makeatletter
\@starttoc{lof}% Print List of Figures
\makeatother

\end{document}
